\section{Conclusion} \label{sec:conclusion}

This paper evaluated three plagiarism detection approaches on a challenging dataset with realistic paraphrasing: TF-IDF baseline, BiLSTM deep learning, and BERT transformers. Contrary to conventional deep learning wisdom, TF-IDF achieves the highest accuracy at 85.51\% with exceptional ROC-AUC of 0.9863, substantially outperforming BiLSTM (59.42\%) and BERT (56.52\%).

Key findings indicate that classical statistical methods (TF-IDF) remain highly effective for plagiarism detection when dealing with realistic paraphrasing patterns (60-85\% word retention with reordering). The performance gap between TF-IDF and deep learning approaches challenges assumptions about semantic understanding automatically improving detection. TF-IDF's 100\% recall ensures all plagiarism cases are caught, with 79.59\% precision minimizing false accusations.

\subsection{Practical Deployment Recommendations}

For academic institutions, TF-IDF-based systems are recommended as the primary screening approach due to superior accuracy (85.51\%), minimal computational requirements (sub-millisecond inference), and perfect recall. Deployment recommendations:

\begin{itemize}
  \item \textbf{Real-time screening} --- TF-IDF only viable option meeting latency requirements
  \item \textbf{Batch processing} --- TF-IDF optimal for comprehensive institutional audits
  \item \textbf{Hybrid approaches} --- TF-IDF for efficient pre-screening with manual review of borderline cases (79-85\% confidence)
  \item \textbf{Ensemble methods} --- Combining TF-IDF with other detectors to further improve precision
\end{itemize}

BiLSTM and BERT show limited practical advantage given TF-IDF's performance superiority and computational efficiency.

\subsection{Future Directions}

Priority research directions include: (1) improved paraphrasing detection through ensemble approaches, (2) multilingual plagiarism detection maintaining TF-IDF's efficiency in multi-language settings, (3) adversarial robustness against intentional evasion techniques, (4) source attribution to identify original content sources, and (5) integration with citation management systems for comprehensive academic integrity platforms.
